% Generated draft: requires mathematical and source review.
\documentclass[11pt]{article}
\usepackage[margin=25mm]{geometry}
\usepackage{fontspec}
\setmainfont{Latin Modern Roman}
\usepackage{amsmath,amssymb,mathtools,unicode-math}
\setmathfont{Latin Modern Math}
\usepackage{enumitem,needspace,xurl,hyperref}
\setlength{\parindent}{0pt}
\setlength{\parskip}{5pt}
\setlength{\emergencystretch}{3em}
\newcommand{\N}{\mathbb{N}}
\newcommand{\Z}{\mathbb{Z}}
\newcommand{\Q}{\mathbb{Q}}
\newcommand{\R}{\mathbb{R}}
\newcommand{\C}{\mathbb{C}}
\newcommand{\F}{\mathbb{F}}
\newcommand{\A}{\mathbb{A}}
\newcommand{\PP}{\mathbb{P}}
\newcommand{\E}{\mathbb{E}}
\newcommand{\eps}{\varepsilon}
\newcommand{\dd}{\,\mathrm{d}}
\begin{document}
\section*{UnsolvedMath ID: 700727}
\textbf{Source problem code:} catalogue entry
\subsection{Definitions and mathematical statement}
A finite group \(G\) is quasisimple if \(G=[G,G]\) and \(G/Z(G)\) is a nonabelian simple group, where \(Z(G)\) is its centre. A complex representation is a homomorphism \(\rho:G\to\operatorname{GL}(V)\), where \(V\) is a nonzero finite-dimensional complex vector space. It is irreducible if its only \(G\)-invariant complex subspaces are \(0\) and \(V\).

To make reduction modulo a prime precise, adopt the following convention. Choose a number field \(K\subseteq\C\) and a \(K\)-representation \(V_K\) with \(\C\otimes_K V_K\cong V\). Write \(\mathcal O_K\) for the ring of integers of \(K\). For each rational prime \(q\), each prime ideal \(\mathfrak p\subset\mathcal O_K\) above \(q\), and each \(G\)-invariant full lattice \(L\subset V_K\) over the local ring \(\mathcal O_{K,\mathfrak p}\), form
\[
\overline V_{\mathfrak p,L}
=\overline{\F}_q\otimes_{\mathcal O_{K,\mathfrak p}/\mathfrak p\mathcal O_{K,\mathfrak p}}
\bigl(L/\mathfrak p L\bigr),
\]
using an embedding of the residue field into \(\overline{\F}_q\). Here a full lattice is a free \(\mathcal O_{K,\mathfrak p}\)-module of rank \(\dim_K V_K\) whose \(K\)-span is \(V_K\). Interpret ``remains irreducible after reduction modulo any prime'' as requiring \(\overline V_{\mathfrak p,L}\) to be irreducible for every such choice of realization, prime, and invariant lattice. In particular, primes dividing \(|G|\) are included. The source does not explicitly state the field or lattice convention; the interpretation above requires review.

\textbf{Problem (catalogue entry; permanent ID 700727).} Classify, up to isomorphism of groups and equivalence of representations, all pairs \((G,\rho)\) in which \(G\) is a finite quasisimple group and \(\rho\) is an irreducible complex representation that remains irreducible modulo every rational prime in the sense just specified. No faithfulness or nontriviality hypothesis is supplied.

The source singles out representations of degree
\[
\frac{p^k-1}{2}
\quad\text{of}\quad \operatorname{Sp}(2k,p),
\]
where \(p\) is an odd prime and \(k\in\N\), as an important example. Here \(\operatorname{Sp}(2k,p)\) denotes the group preserving a nondegenerate alternating bilinear form on \(\F_p^{2k}\). The source does not identify the representations within this degree or state parameter exceptions. Its convention for \(\N\) is also unspecified; \(k=0\) would give degree zero and hence no representation under the nonzero convention above. Any example within the requested classification must additionally satisfy the quasisimplicity hypothesis.
\subsection{Short English statement}
Classify the irreducible complex representations of finite quasisimple groups whose reductions remain irreducible in every prime characteristic. The draft makes the reduction convention explicit and retains the source's symplectic example with its unresolved parameter and representation qualifications. The supplied registry labels the problem open; this status has not been independently verified.
\subsection{Sources}
\begin{itemize}
\item Supplied local UnsolvedMath registry record, permanent ID \(700727\), exact source problem code \texttt{catalogue entry}, title \emph{Kourovka Notebook Problem 11.34}. The supplied status is \texttt{open}.
\item \emph{The Kourovka Notebook (Unsolved Problems in Group Theory)}, original problem 11.34, 11th issue (1990); the supplied record also cites the 21st issue edition (2026), \texttt{arXiv:1401.0300v48}. The collection description names E. I. Khukhro and V. D. Mazurov as editors. These edition details are supplied attribution, not independently verified.
\item Proposed in 1990 by A. E. Zalesski; the supplied proposer field is truncated and reads \texttt{A. E. Zalesski\textbackslash u\{\textbackslash i}. The intended final diacritic is unresolved.
\item P. H. Tiep and A. E. Zalesski, \emph{Proc. London Math. Soc.} \textbf{84} (2002), 439--472. Cited by the source for partial results; no article title is supplied.
\item P. H. Tiep and A. E. Zalesski, \emph{Proc. Amer. Math. Soc.} \textbf{130} (2002), 3177--3184. Cited by the source for partial results; no article title is supplied.
\item C. Parker and M. van Beek, preprint (2024), \url{https://arxiv.org/abs/2411.16379}. Listed in the source's discussion of partial results; no title or precise result is supplied.
\end{itemize}
\end{document}
